\subsection{Architect-Specified Strategies for NaturalPlan}

\paragraph{Trip Planning.}
The Architect models trip planning as a constrained path-finding problem over
cities.  A city is treated as a node, a direct flight is treated as a directed
edge, and a complete trip is an ordered path that visits every required city.
The important timing idea is that each city visit has an inclusive day
interval, and the flight day is shared between two consecutive cities.  For
example, if a visit to city \(A\) ends on day \(d\), then the next city starts
on the same day \(d\).  Fixed events are handled as interval constraints: if
the prompt says that the traveler must be in a city between two dates, then
that whole event window must lie inside that city's visit interval.

The implemented code follows this design by first parsing cities, stay
durations, total trip length, flight edges, and fixed event windows.  It then
checks whether the total duration is arithmetically consistent with the total
number of trip days.  After that, it runs a deterministic depth-first search
over possible flight paths.  During search, it rejects a partial path if the
current city interval violates an event constraint, exceeds the total trip
length, or makes some remaining city impossible to schedule later.  For small
cases, the code also has a permutation fallback over city orders.

This strategy is not a hallucination.  It is a standard graph-search and
constraint-satisfaction formulation, similar to searching for a Hamiltonian
path with temporal constraints.  The strategy is reliable for the NaturalPlan
trip instances supported by the parser because the final route is checked for
city coverage, valid flights, correct durations, correct total length, and
event-window satisfaction.

\paragraph{Meeting Planning.}
The Architect models meeting planning as a time-window routing problem.  The
state of the planner consists of the current location, the current clock time,
and the set of friends already met.  Each friend has a location, an
availability window, and a minimum required meeting duration.  A valid next
step must account for travel time from the current location, possible waiting
until the friend becomes available, and enough remaining time to complete the
meeting before the friend's window closes.

The implemented code parses the start location and time, each friend's
location, availability window, required duration, and the travel-time matrix.
It then uses memoized dynamic search over states.  For each state, the solver
tries every friend who has not yet been met.  It computes arrival time,
meeting start time, and meeting end time.  If the meeting cannot finish inside
the availability window, that option is discarded.  Otherwise, the solver
recursively continues from the friend's location and the new time.  The chosen
schedule maximizes the number of friends met, with deterministic tie-breaking
for earlier completion, less waiting, and stable input order.

This is an existing algorithmic idea, not a hallucinated one.  It is a dynamic
programming search over subsets, closely related to exact approaches for
small routing, scheduling, and orienteering problems with time windows.  The
strategy is reliable for validity because the code verifies travel times,
arrival times, time-window constraints, meeting durations, and that no friend
is repeated.  Exact textual match may still differ when multiple optimal
meeting orders are valid.

\paragraph{Calendar Scheduling.}
The Architect models calendar scheduling as selecting one feasible meeting
slot from a finite calendar grid.  The problem is represented using the
participants, allowed days, work-hour bounds, meeting duration, busy intervals,
and optional preference constraints such as blocked days or lower/upper time
bounds.  A candidate slot is valid only if it fits within work hours, has the
required duration, and does not overlap with any participant's busy intervals
or parsed preferences.

The implemented code follows this representation directly.  It parses the
task header to obtain participants, duration, work hours, and candidate days.
It then parses each participant's busy intervals and merges overlapping
intervals.  Preference statements are converted into blocked days,
not-before constraints, or not-after constraints.  The solver scans candidate
start times on a 30-minute grid in chronological order.  For each candidate
slot, it checks all participant constraints.  The first slot that passes these
checks is returned as the meeting time.

This is also not a hallucinated strategy.  It is a standard interval
scheduling and constraint-enumeration method.  Since the candidate time grid
is finite, the algorithm can simply test possible slots until it finds a
valid one.  The strategy is reliable for the NaturalPlan calendar format
because the returned slot is verified against duration, work hours, busy
intervals, blocked days, and parsed time preferences.

\paragraph{Overall assessment.}
The Architect's choices are algorithmically grounded across all three
domains.  Trip planning uses constrained graph search, meeting planning uses
dynamic search over time-window routing states, and calendar scheduling uses
interval-based slot enumeration.  These are recognizable planning and
scheduling methods, not invented or hallucinated algorithms.  The main
limitation is that reliability depends on the prompt being parsed into the
intended structured constraints.
